\section{결과}
\label{sec:results}

\rev{\ref{sec:method:samples}절의 점검에서, 204경기(코호트마다 17경기)의 1,293개 경계 모두에서 다시 만든 입력이 원래 입력과 같았고, 절단을 뺀 빌더는 시험한 3경기의 21개 경계 모두에서 누수로 검출되었으며, 817,766개 표본 가운데 절단 시각 이전 기록으로 살아 있는 팀원이 없던 표본은 2개(시험 표본에는 없음)였다.}
모든 차이는 12개 코호트의 확률효과 통합 추정치와 95\% 신뢰구간이며, 괄호 안의 분수는 통합 추정치와 부호가 같은 코호트의 수이다(12/12는 모든 코호트에서 방향이 같다는 뜻이다).
AUC 차이 0.01은 탈락 팀에 생존 팀보다 높은 위험을 줄 확률이 1\%p 높다는 뜻이며, \rev{개선과 실질적 동등은 \ref{sec:method:evaluation}절의 기준으로 판정한다.}

\begin{figure}[t]
\centering
\includegraphics[width=\columnwidth]{fig_rq1_effects.pdf}
\caption{단계 구간별 서식지 맥락의 효과}
\label{fig:rq1}
\end{figure}

\begin{table*}[t]
\centering
\caption{설계별 예측 성능과 원래 과제 성능}
\label{tab:main}
\footnotesize
\setlength{\tabcolsep}{4pt}
\renewcommand{\arraystretch}{1.15}
\begin{tabular}{lcccccccc}
\hline
 & \multicolumn{6}{c}{본 과제: 다음 단계 탈락} & \multicolumn{2}{c}{KRAFTON의 원래 과제} \\
\cline{2-7}\cline{8-9}
설계 & AUC & 단계 안 AUC & 경기 안 AUC & 로그 손실 & Brier & 단계별 보정 오차 & 우승 팀 정확도 & C-index \\
\hline
본 연구, GBDT & 0.778 & 0.752 & \textbf{0.757} & \textbf{0.489} & \textbf{0.160} & \textbf{0.013} & 0.351 & 0.654 \\
본 연구, PTR-Former & \textbf{0.779} & \textbf{0.757} & \textbf{0.757} & 0.490 & \textbf{0.160} & 0.017 & 0.386 & 0.663 \\
KRAFTON 적용판 & 0.741 & 0.741 & 0.748 & 0.518 & 0.171 & 0.064 & \textbf{0.417} & \textbf{0.665} \\
KRAFTON + 단계 기저위험 & 0.763 & 0.741 & 0.748 & 0.504 & 0.165 & 0.017 & \textbf{0.417} & \textbf{0.665} \\
Sagong 적용판(XGBoost) & 0.686 & 0.632 & 0.645 & 0.562 & 0.188 & 0.016 & --- & --- \\
\hline
\multicolumn{9}{l}{\emph{같은 KRAFTON 구조에서 학습 목적만 바꾼 경우}} \\
Cox 목적 & 0.742 & 0.731 & 0.745 & 0.517 & 0.171 & 0.056 & \textbf{0.387} & \textbf{0.661} \\
다음 단계 교차엔트로피 목적 & \textbf{0.776} & \textbf{0.755} & \textbf{0.753} & \textbf{0.491} & \textbf{0.161} & \textbf{0.016} & 0.362 & 0.653 \\
\hline
\multicolumn{9}{l}{\emph{\rev{두 입력을 합친 경우}}} \\
GBDT, 두 입력 결합 & 0.793 & 0.769 & 0.774 & 0.475 & 0.155 & 0.013 & 0.369 & 0.662 \\
\hline
\end{tabular}
\\[2pt]
\parbox{\textwidth}{\scriptsize 12개 코호트 평균\rev{(본문의 차이는 확률효과 통합 추정치라 표의 평균 차이와 조금 다를 수 있다)}. 굵은 글씨는 각 블록에서 가장 좋은 값이다. AUC와 정확도, C-index는 클수록, 로그 손실, Brier, 보정 오차는 작을수록 좋다. 우승 팀 정확도와 C-index는 KRAFTON 구현이 정의한 지표이며, 본 연구의 PTR-Former는 순위 출력을 사용하였다.}
\end{table*}

\begin{figure*}[t]
\centering
\includegraphics[width=\textwidth]{fig_contrast.pdf}
\caption{단계별 보정과 설계 요인의 분해\rev{((b), (c)의 막대는 12개 코호트 평균, 오차 막대는 코호트 간 95\% 구간)}}
\label{fig:contrast}
\end{figure*}

%======================================================================
\subsection{RQ1: 서식지 맥락과 다음 단계 탈락 위험}
\label{sec:results:rq1}

\rev{그림~\ref{fig:rq1}의 계수는 요인이 한 표준편차 늘 때 탈락 위험의 로그 오즈가 바뀌는 양이며, 지수를 취하면 탈락 오즈가 몇 배가 되는지를 뜻한다.}
모든 단계를 합친 모형에서 주변 경쟁 밀도는 $+0.293$ $[0.242, 0.343]$(12/12)로, 주변의 적 팀이 한 표준편차 많으면 탈락 오즈가 약 1.34배가 되었다.
이동 부족분은 $+0.067$ $[0.051, 0.083]$(12/12, 약 1.07배), 남은 팀원 수는 $-0.632$ $[-0.713, -0.551]$(11/11, 약 0.53배)였다.
\rev{세 요인 모두 평가 가능한 모든 코호트에서 방향이 같았으므로, 경쟁과 이동 부담이 클수록, 남은 팀원이 적을수록 위험이 높다고 해석된다.}
\rev{반면 구역 밖 깊이의 통합 추정치는 $-0.031$ $[-0.056, -0.007]$(8/12)로 작았고 단계 구간마다 부호가 바뀌어, 일관된 연관으로 보기 어려웠다. 작은 음의 부호는 이동 부족분이 같을 때 구역 밖으로 더 떨어져 있다고 위험이 더 커지지는 않았다는 뜻일 뿐이다(두 변수의 분산팽창계수(VIF)는 2.9 이하).}
\rev{강건성 재적합에서 경쟁 밀도와 남은 팀원 수는 AI 팀 지표를 더해도($+0.304$, $-0.613$), 반경을 100\,m나 500\,m로 바꿔도($+0.242$, $+0.194$), 경쟁전 경기만 써도(6개 코호트; $+0.281$, $-0.715$) 모든 코호트에서 부호가 유지되었고, 구역 안쪽 깊이도 음수로 남았다. 반면 이동 부족분은 AI 팀 지표를 더하면 $+0.017$ $[0.003, 0.031]$(8/12), 경쟁전 경기만 쓰면 $+0.028$ $[-0.012, 0.069]$로 줄었다. AI 팀이 이동 부족분도 크고 탈락 위험도 높기 때문이다.}

\rev{연관의 크기는 단계에 따라 달랐다(그림~\ref{fig:rq1}은 구간 안 표준편차 단위).}
구역 안쪽 깊이의 계수는 초반 $-0.028$(신뢰구간이 0을 포함, 7/12)에서 중반 $-0.289$, 후반 $-0.410$(모두 12/12)으로 커졌는데, \rev{이는 구역이 작아질수록 경계에서 깊이 들어가 있는 팀의 위험이 더 낮았다는 뜻이다.}
\rev{이동 부족분은 초반과 후반에 양수였지만($+0.066$, 11/12; $+0.097$, 9/12) 중반에는 0에 가까워($+0.006$, 6/12), 작은 효과가 단계에 따라 불안정하였다.}
경쟁 밀도의 계수는 초반 $+0.376$에서 후반 $+0.229$로 작아졌고, 구간 간 차이는 우연으로 보기 어려웠다($p=5.9\times10^{-5}$).
\rev{공통 단위(학습 집합 표준편차)로 다시 추정해도 안쪽 깊이는 $-0.02$, $-0.87$, $-6.13$으로 커지고 경쟁 밀도는 $+0.58$에서 $+0.20$으로 작아져, 구간 간 차이가 유지되었다(각각 $p=5.6\times10^{-11}$, $1.4\times10^{-6}$).}
\rev{확률 단위의 평균 한계효과로 보면,} 주변의 적 팀 하나당 다음 단계 탈락확률은 초반 \rev{9.0}\%p, 중반 \rev{4.2}\%p, 후반 3.8\%p 높았고, 팀원 한 명당 각각 \rev{7.4}\%p, \rev{11.4}\%p, \rev{13.6}\%p 낮았다.

\rev{이 요인들은 적은 수의 변수로도 위험의 상당 부분을 설명한다.
GBDT가 우연 수준(AUC 0.5)을 넘어 얻은 판별력 가운데 단계 절편만 쓴 모형은 40.4\% $[38.5, 42.5]$를, 네 변수를 더한 간이 모형은\rev{(두 깊이 항을 뺀 모형으로, 바깥 깊이와 필요 속도의 상관은 $r=0.89$--$0.95$이며 두 항을 다시 넣으면 통합 AUC가 0.002, 단계 안 AUC가 0.005 오른다)} 78.8\% $[77.2, 80.3]$를 재현하였고, 단계 수준의 몫이 빠지는 단계 안 AUC로 재도 간이 모형의 몫은 70.6\% $[68.1, 73.5]$였다.}

%======================================================================
\subsection{RQ2: 기존 설계와의 비교}
\label{sec:results:rq2}

표~\ref{tab:main}에서 본 연구의 설계는 AUC 0.778(GBDT)과 0.779(PTR-Former)로, KRAFTON 적용판(0.741), 단계 기저위험을 결합한 KRAFTON(0.763), Sagong 적용판(0.686)보다 높았다.
PTR-Former와 기저위험을 결합한 KRAFTON의 차이는 $+0.016$ $[0.014, 0.019]$로, 탈락 팀과 생존 팀의 쌍을 올바르게 줄 세울 확률이 1.6\%p 높고 12/12 코호트에서 같은 방향이어서 개선으로 판정된다\rev{(같은 기준에 대한 GBDT의 차이는 $+0.013$ $[0.010, 0.015]$, 12/12)}.
Sagong 적용판과의 차이는 $+0.089$ $[0.082, 0.096]$(12/12)로 더 컸다.

그림~\ref{fig:contrast}(a)는 \rev{시험 표본에서} 단계마다 실제 탈락률(검은 선)과 각 설계의 평균 예측확률을 겹쳐 보이며, 두 선이 겹칠수록 그 설계가 단계별 위험 수준을 바르게 나타낸다는 뜻이다.
실제 탈락률은 단계 1의 19.5\%에서 단계 7의 50.6\%로 오르지만, KRAFTON 적용판의 예측은 26.4\%에서 29.8\%로 거의 평평하여 초반에는 위험을 과대평가하고 단계 7에서는 과소평가한다.
단계별 보정 오차는 KRAFTON이 0.064(단계별 평균 예측이 실제와 평균 6.4\%p 어긋남), 본 연구의 GBDT가 0.013이었고, KRAFTON에 기저위험을 결합하면 0.017로 줄었다.
\rev{이 지표는 단계 평균만 비교하므로 단계 절편을 가진 설계라면 모두 낮으며, Sagong 적용판도 0.016이다.
단계 안 ECE로 보면, 본 연구의 GBDT 0.039와 PTR-Former \rev{0.042}가 기저위험을 결합한 KRAFTON(0.046)과 KRAFTON 적용판(0.077)보다 작았다.}

\rev{따라서 설계의 차이는 단계 수준을 뺀 비교로 확인해야 한다.
\rev{기저위험을 결합한 KRAFTON을 기준으로 하면, 같은 단계의 표본끼리만 비교한 단계 안 AUC의 개선량은 PTR-Former가 $+0.017$ $[0.013, 0.021]$(10/12), GBDT가 $+0.007$ $[0.003, 0.011]$(11/12)였고, 같은 경기의 같은 순간에 살아 있던 팀끼리만 비교한 경기 안 AUC의 개선량은 각각 $+0.010$과 $+0.009$(모두 11/12)였다. PTR-Former와 GBDT의 직접 비교는 \ref{sec:results:rq3}절에서 다룬다.}
두 지표는 단계마다 예측을 함께 올리거나 내리는 조정에 영향을 받지 않으므로, 이 차이는 기저위험이나 확률 보정 방식에서 생긴 것이 아니다.}

표~\ref{tab:main}의 오른쪽 두 열은 KRAFTON의 원래 과제 성능으로, 우승 팀 정확도는 \rev{각 절단 시각에} 살아 있는 팀 가운데 최종 우승 팀을 1위로 맞힌 비율(무작위 0.182)이고 C-index는 탈락 순서를 맞힌 팀 쌍의 비율이다.
이 과제에서는 KRAFTON 적용판이 0.417과 0.665로 가장 높았고 본 연구의 GBDT는 0.351과 0.654였으며, 같은 KRAFTON 구조에서도 Cox 목적은 원래 과제에서, 다음 단계 교차엔트로피 목적은 본 과제에서 더 높았다.
\rev{Sagong 적용판도 분석 단위가 같은 솔로 코호트의 개인 생존 분류에서 모두 생존으로 예측하는 기준보다 정확도가 1.8\%p $[0.1, 3.3]$ 높아, 원 논문의 표에서 계산한 개선 폭(3.3\%p)보다 작지만 같은 방향이었다.
\rev{같은 특징 코드와 학습기로 사후 구성(단계 시작의 표본, 단계 전체의 지표, 그 단계 끝까지의 생존을 라벨로)을 쓰면 솔로 AUC가 0.653에서 0.833으로 올랐다. 원 논문은 이 시점을 명시하지 않는다.}
따라서 두 적용판은 원래 과제에 맞게 학습되었다.}

%======================================================================
\subsection{RQ3: 설계 요인의 분해}
\label{sec:results:rq3}

\begin{table}[t]
\centering
\caption{\rev{지표별 분해 단계의 AUC 변화}}
\label{tab:steps}
\footnotesize
\setlength{\tabcolsep}{3.5pt}
\begin{tabular}{lrrr}
\hline
단계 & 통합 & 단계 안 & 경기 안 \\
\hline
(a) Sagong 입력 $\to$ 본 연구 입력 & \textbf{+0.076} & \textbf{+0.103} & \textbf{+0.097} \\
(b) KRAFTON 입력 $\to$ 본 연구 입력 & +0.001 & $-$0.002 & $-$0.004 \\
(c1) Cox 점수 + 단계 기저위험 & \textbf{+0.016} & 0$^\dagger$ & 0$^\dagger$ \\
(c2) Cox $\to$ 다음 단계 교차엔트로피 & \textbf{+0.018} & \textbf{+0.022} & \textbf{+0.009} \\
(d) 명시적 단계 절편 & $-$0.001 & $-$0.001 & +0.000 \\
(e) KRAFTON $\to$ 팀 축 PTR-Former & +0.001 & +0.002 & +0.002 \\
(f) + 단계 축·관계 편향 & +0.002 & +0.002 & +0.001 \\
(g) + KRAFTON 입력(합집합) & \textbf{+0.014} & \textbf{+0.018} & \textbf{+0.016} \\
\hline
\end{tabular}
\\[2pt]
\parbox{\columnwidth}{\scriptsize 확률효과 통합 추정치. 굵은 글씨는 95\% 신뢰구간이 0을 포함하지 않고 크기가 0.005 이상인 변화이다. (a), (b), (g)는 GBDT, (c1)--(d)는 KRAFTON 구조에서 잰다. $^\dagger$단계마다 단조 변환이므로 정의상 0이다.}
\end{table}

\rev{표~\ref{tab:steps}는 각 단계의 통합 변화량을 세 지표로 보이고,} 그림~\ref{fig:contrast}(b)와 (c)의 막대 하나는 요인 하나를 바꿀 때 \rev{12개 코호트 평균} AUC가 변한 양이며, 진한 막대가 \rev{핵심 단계이다(그림은 평균이라 막대를 더하면 끝값과 같다)}.
Sagong 경로(그림~\ref{fig:contrast}(b))에서는 같은 GBDT에 서식지 맥락 블록만 더해도 \rev{평균} AUC가 \rev{0.699에서 0.051} 올랐고, 나머지 입력을 본 연구의 나머지 팀 특징으로 바꾸면 \rev{0.027}이 더해졌다.
\rev{전체 차이의 약 65\%가 서식지 맥락 블록에서 나왔고, 거꾸로 본 연구의 전체 입력에서 이 블록만 빼도 GBDT의 통합 AUC가 0.012(단계 안 0.013, 모두 12/12) 낮아져, 이 정보는 나머지 특징으로 대체되지 않았다.}

\rev{KRAFTON 경로(그림~\ref{fig:contrast}(c))에서 입력과 학습 표본을 본 연구에 맞추는 단계는 평균 AUC를 바꾸지 않았지만(0.000) 단계 안 AUC는 통합 추정치로 0.010 $[0.005, 0.015]$ 낮췄다.}
반면 Cox 점수에 기저위험을 결합하면 \rev{0.017}, 학습 목적을 다음 단계 교차엔트로피로 바꾸면 \rev{0.018}이 더해졌고(모두 12/12), 목적 교체의 이득은 단계 안 AUC($+0.022$, \rev{12/12})와 경기 안 AUC($+0.009$, \rev{11/12})에서도 남았다\rev{(표~\ref{tab:steps})}.
명시적 단계 절편($-0.001$), 아키텍처(\rev{0.000}), 단계 축과 관계 편향(\rev{$+0.003$})은 \rev{변화가 작았고 통합 AUC로} 모두 실질적 동등이었다.
따라서 \rev{이 경로에서 순위 손실 없는 PTR-Former까지의} 전체 차이 0.037 가운데 \rev{약 94\%(0.035)}가 기저위험과 학습 목적에서 나왔다\rev{(두 단계는 같은 수준 정보를 일부 공유하므로 각 몫보다 합이 안정적이다)\footnote{순서를 바꾸어 학습 목적부터 바꾸면 그 단계만으로 $+0.034$ $[0.029, 0.040]$(12/12), 곧 두 단계를 합친 만큼이 더해진다. 교차엔트로피가 이미 단계별 수준을 맞히기 때문이다(보정 오차 0.016, 표~\ref{tab:main}). 두 단계의 몫은 순서에 따라 달라지지만 합은 달라지지 않는다. 학습 목적의 단계 안·경기 안 이득은 순서와 무관한데, 기저위험이 각 단계 안에서 단조 변환이기 때문이다.}}.

\rev{구조의 효과는 지표에 따라 달랐다.}
사전 지정한 다섯 비교(네 개는 통합 AUC, 하나는 순위 상관)는 모두 실질적 동등이었고 REML과 Hartung--Knapp 방식으로도 판정이 같았으며, 팀 축의 효과는 $+0.0027$ $[0.0015, 0.0038]$(Holm $p<0.001$)로 0.3\%p에 그쳤고\rev{ 같은 경기·같은 단계 안의 순위 상관에도 의미 있는 변화가 없었다($+0.0002$). 다만 이 판정은 SESOI에 달려 있어, SESOI가 0.0036보다 작으면 팀 축 효과는 판정 불가, 0.0026 이하이면 개선이 된다.
그러나 같은 단계 안의 비교에서는 팀 축이 $+0.010$ $[0.008, 0.012]$(11/12), PTR-Former가 GBDT보다 $+0.009$ $[0.006, 0.012]$(11/12) 높았다.
같은 경기 안의 비교에서는 아키텍처($+0.002$)와 단계 축·관계 편향($+0.001$)이 0과 구분되지 않았고 두 모델의 경기 안 AUC도 같았으므로(0.757), 팀 축은 같은 단계라도 경기마다 다른 위험 수준을 담는 데 기여한 것으로 보인다. KRAFTON 구조에도 팀 간 attention이 있으므로 (e)의 변화는 작았다.}
\rev{현재 학습 설정에서 최대 학습 횟수를 두 배로 늘린 실험은 구조의 기여가 작다는 결과를 뒤집지 않았고(10개 실행 중 9개에서 AUC $-0.0001$에서 $-0.0034$), 순위 손실의 효과도 통합 AUC로 0.001 이하였다.}
\rev{또한 가장 큰 두 코호트에서 초기화 세 개에 걸친 AUC의 표준편차는 모델별 0.0002--0.0018, 비교별 0.0003--0.0022로 모두 SESOI(0.005)보다 작아, 검증한 두 코호트에서는 동등 판정이 초기화에 따라 바뀌지 않았다.}
\rev{패치 40.2와 41.1로 학습하고 이후 패치 41.2로 시험하면, 같은 시험 행에서 GBDT의 AUC는 $0.0021$ $[0.0008, 0.0035]$(11/12), 단계 안 AUC는 $0.0039$ $[0.0003, 0.0075]$ 낮아졌다. 학습 경기 수를 같게 맞춘 무작위 분할은 $0.0003$ $[-0.0015, 0.0021]$ 낮아지는 데 그쳤다. 사전 규칙으로 보면 패치 변화는 통합 AUC에서 실질적 동등, 단계 안 AUC에서는 판정 불가이다.}
끝으로 본 연구의 입력과 재구성한 KRAFTON 입력은 같은 GBDT에서 실질적으로 같았지만($+0.001$ $[-0.002, 0.004]$), 둘을 합치면 $+0.014$ $[0.013, 0.016]$(12/12) 높아져\rev{, 합친 GBDT의 AUC는 0.793이었다(표~\ref{tab:main})}.
